\chapter[Sermon 71. The Sixth Angel Sheds His Vial.]{The Sixth Angel Sheds His Vial.}
\label{sermon:071}
\markright{Sermon 71. The Sixth Angel Sheds His Vial.}

And the sixth angel pours out his vial upon the great river Euphrates, and the water dried up, that the way of these kings of the East should be prepared. And I saw three unclean spirits, like frogs, come out of the mouth of the Dragon, and out of the mouth of the beast, and out of the mouth of the false prophet. For they are the spirits of the devils working miracles, to go out to the kings of the earth, and of the whole world, to gather them to the battle of the great day of God almighty. Behold, I come as a thief. Happy is he that watches, and keeps his garments, lest he be found naked, and men see his shame. And he gathered them together in a place called in the Hebrew tongue Armageddon.

\index{Daniel (Book of)}\index{Rome}The sixth angel pours out his vial on the great river Euphrates; the end of this shedding is that the way might be opened for the kings of the East, that is, that Babylon might be taken. This plague chiefly pertains to Rome and the Roman church. The speech has an allegory, or a secret comparison and allusion to old Babylon. We read in the fifth chapter of Daniel that Babylon was taken the same night wherein Balthazar, king thereof, had prepared a sumptuous banquet and looked for nothing less than destruction. Herodotus and Xenophon report how the kings of the East, Darius Priscus, who is also called Medes, and Cyrus of Persia besieged the city round about; but where there was no hope to win it, Cyrus let out the Euphrates by ditches, so that the army might wade over the river; and so was the city laid open and taken on the side where it was fenced with the river. Euphrates therefore fortified Babylon and brought to it many other commodities and pleasures. Here are signified, by Euphrates, riches, defense, pleasures, commodities, tributes, and customs, which the Roman churches call sacred or of the holy church. These commodities and pleasures are diminished by the kings of the East, by true Christians, whom the Scripture calls the Kings and Priests, and derived and put to another use. Wherefore the power of the Roman church begins to decay, to the intent that at length it may be taken and abolished by the Lord Christ himself. Doubtless the true Christians understand, believe, and profess that Christ alone is the Savior, neither that there is salvation in any other. And that this is given freely; that they be mad and commit Simonry and sacrilege, which in this case practice and make merchandise. Read the fifty-fifth chapter of Isaiah, and the eighth of the Acts, finally the first and second chapter of Paul to the Colossians, wherein most diligently is declared, illumined, and set forth that by Christ alone we are absolved and in him alone have all fullness. And what time the common people do hear this, to wit that by those Roman trifles, fairs of pardons, and other crafty jugglings they are deceived and robbed of their substance, they shut by and by and make fast their chests, their purses, their cellars, and garners. And so dries up the river of wealth and pleasure; it dries up also when the godly deny to give other customs, as tithes, palles [?], first fruits, and such other like things. So is the way prepared for the kings of the East, so begins Rome, the second Babylon, to be taken and come to naught.

\index{Antichrist}\index{Papacy}It follows furthermore, how Antichrist will fight against the faithful and godly laboring to dry up Euphrates, for the maintenance and increase of his kingdom: and where he might briefly have said, he shall send forth ambassadors unto all kings and princes, to stir them up against the gospelers, for the defense of the privileges, rights and revenues of the See of Rome, he had rather most diligently describe those ambassadors and show their destruction. It forces very much to have known the Popes legates. For they are marvelous pestilent to the church of God: for we have not only experience of it at this day, but also by the reading of all histories, that great evils and all calamities in a manner have been brought into the church, and are also at this day, through the instigations of those legates. I touched a little before what mischief Cardinal Julian Cesarine, the legate of Pope Eugenius, wrought in Germany, Bohemia, Poland and Hungary. What is done in our time, and has been done in our fathers’ memory, it were too long to rehearse. If our elders had by the doctrine of Jesus Christ revealed to the church by John understand and known the nature of the Popes legates, they might easily have eschewed wherewith they have undiscreetly entangled themselves, and suffered great loss and hindrance. I speak nothing here of ambassadors and embassies of kings and commonwealthes uncorrupted.

\index{Primasius}First, he diligently demonstrates the origin of legates, so that we may understand that they are led by a wicked spirit, and that their vocation is not godly, but devilish. He shows a threefold origin, within which, in truth, they may all be reduced to one devilish unity. He says that Primasius, explaining this passage, saw one spirit, and because of the number of the parts of one body, he says three, so that the number of the wicked might be expressed as being led by one devilish spirit. Therefore, the first den from which the legates break out, he calls the dragon’s mouth. The dragon is spoken of in chapter 12; and there is no one who does not understand that it signifies the devil himself. They therefore come forth from the devil. For all the affairs of their embassy consist in lies, schemes, practices, finally, in corrupting the truth and sincerity of the gospel, and also in factions and dissensions, in slaughter and bloodshed. And the devil was from the beginning a liar and a murderer, as the Lord himself says in John 8. And so far, they are from the dragon’s mouth. The same arise also from the mouth of the beast. For they come furnished with the Pope’s authority, legates lateral with full power. Of the beast, I have spoken in chapter 13. Finally, they come out of the mouth of the false prophet. The true prophet and pastor, supreme and only of the universal church, is Christ, the Son of God. Antichrist is that false prophet and chief seducer of the whole world, as is said in chapter 13. Therefore, the legates come, sent from the Pope, who has put into their mouths words, instructions, or commissions that they should speak those things which are of false prophecy. However, explaining himself more plainly, he declares of what sort the legates shall be: namely, three unclean spirits. An unclean spirit is everywhere in Scripture called the devil or Satan, truly in nature and effect. For as the Spirit of God is called holy, so this is, contrarily, unclean. For he himself is, by nature, or rather by his own corruption and revolt from God, impure, and the author of all impurity and uncleanness. He signifies therefore that those legates shall be men of a devilish uncleanness. And indeed, if you add to this the lives, manners, and conversations of those lateral legates, and of their families, you will find, in effect, nothing else but extreme uncleanness, filthiness, and beastliness, monstrous lust, whoredom and adultery, and detestable fornications, wonderful surfeiting, bloody schemes and counsels. Therefore, the thing itself speaks: and the things that the legates do everywhere are a commentary on this passage. And where three unclean spirits are reckoned, some explain it of divines, lawyers, and religious persons, such as monks and friars, from which three sorts the Pope’s ambassadors are mostly chosen. I understand simply by the third number that those legates shall be most furnished with all hostile authority, and that they shall all agree well among themselves, and all help one another: so that whatever one seems to lack, another may supply. Solomon in Ecclesiasticus says, “A threefold rope or line will not easily break.”

But now, so that no one should lack understanding, he presents, through a parable, what kind of men these legates will be—truly frogs of the marsh or fen, and insistent, tedious criers, foul and filthy. And he does not say that they are frogs in reality, but like frogs. For just as frogs, by their insistent crying, are most tedious and troublesome, and those of the marsh are also filthy, so these legates love earthly things and filthiness. And by their complaints, accusations, provocations, writings, and arguments, they are altogether frog-like and fen-like, hateful both to God and men. They are nothing ashamed; if they are interrupted a little, they immediately return to their old song, [?persisting]. For there is no other tune with them but [?that]. Primasius reasons greatly about frogs. Among other things, it is fitting for false prophets, like frogs crying in the night, to make a damnable noise by barking out errors. For frogs, because of their location, appearance, and troublesome noise, are so hateful that the Devil, with his works, is known to be abominable to the truth, and with just fire rightly condemned, and so forth. Thus he says. And just as the frogs of Egypt were raised from the dust by the devilish art of the magicians, and cried out against God’s truth, calling back the people of God through Moses and Aaron to true liberty and worship of God, so too the Pope’s legates harass with talk the preaching of the gospel, the free deliverance, the Christian liberty, and true service of God. And just as the frogs double and repeat even to the point of making one weary to hear, so these fen-like beasts of Rome always have in their mouths, “the most holy See,” “the most holy father,” “the holy church of Rome.” “The holy church of Rome does not err,” “the holy church of Rome must be obeyed.” “He who will not obey her is a heretic and a schismatic.” These things they will reiterate and repeat many times and often, to all people, and in all and every cause, their one and the same song [?sounding].

The Lord adds through John, and thus declares it even more clearly: for they are spirits of devils working miracles, \&c. “Daemon” (which is here used in Greek for devil) has its name from sundry knowledge and skillfulness of things, and seems to be in a manner indifferent, although it is commonly put for the Devil. Nevertheless, for a difference, they are called Eudaimones and Cacodaimones, as if good and evil workers. For the Greeks say that “daemon” is called from “daina” that is knowing or skillful. For “demiurgos” is called an expert Artificer. The Lord therefore signifies that the Popes’ legates shall be spirits of devils, that is to say, spiritual fathers (but endowed with the spirit of Satan), wise men or skillful, crafty workers to bring their matters to pass. And therefore he adds, working wonders. Whereby he seems to allude to the magicians of Egypt, who also wrought miracles and detained King Pharaoh in lies against the truth. Paul moreover in 2 Timothy 3 compares the wise men and ministers of Antichrist to the magicies of Egypt. And it is well known that the legates do everywhere boast of miracles which have been done in their church and religion, and so keep still the hearts of kings and princes in papist errors. Of miracles speaks Paul in 2 Thessalonians 2. And I have said something hereof in Chapter 13.

\index{John, Saint}\index{Jerome, Saint}Moreover, here is shown the end of all the treatises and counsels of the Popes’ legates: that they might go forth to the kings of the whole earth, to assemble them to battle. And so they shall creep into the courts of all kings and princes. You will surely find the Popes’ legates in a manner in all the kings’ courts. And what do they do? They surround kings and princes. They see that no faithful man be admitted to the kings’ speech, they learn to know all the kings’ counsel, which they write and signify to Rome: and if they dislike anything, that they may invalidate and subvert the same; and that they always repeat that song of theirs, namely [illegible], that is truly obedience, which all men owe to the Holy See. Finally, that they arm kings and princes to defend the church of Rome and destroy heresies. This, I say, is the battle of that same day of the great God almighty, that is to say, which shall be marked by the coming of the Son of God unto judgment, and which shall endure to the coming of Christ unto judgment, when he shall then avenge the blood of his, from the hands of that hideous beast. And he calls the day of judgment, the day of the great God, as does also Jerome in the second letter to Titus. And the day of God almighty: as he who shall then show his omnipotence, and even his divine power, which seems now to the ungodly by reason of his long forbearance to sleep. This necessary and most profitable description Saint John has set here, by the revealing of Jesus Christ, so that we should watch and beware of them.

Hereafter follows a faithful admonition and exhortation to watchfulness, lest we fall asleep and perish with the Antichristians in the cares and pleasures of this world. And he says how that day of the Lord will come suddenly, and when we shall least look for it. For the Lord here repeats that thing which he said also in the Gospel: behold, I come as a thief. These things are read in Matthew 24 and are repeated by the Apostle in 1 Thessalonians 5. And verily, that same sudden coming of the Lord excites the minds of us all and provokes to watchfulness, lest we should at unawares be oppressed. He adds also immediately a profit prepared for them that watch. Happy, says he, is that man that watches. He adds moreover, how the godly should demean themselves in watching. How they must keep their garments, that they be not defiled; and take heed moreover that they walk not naked, lest their filthiness be espied. Touching garments I have spoken most largely in another place of this book. He keeps his garments, that keeps his life and conversation unspotted of worldly filthiness. He does not walk naked, which puts on Christ. But his shame is seen, that sins impudently. But chiefly is their shame seen, whose whoredoms, adulteries and fleshly lusts are known and open to the eyes of all men. And here is the state of them to be lamented who are called spiritual, and rather in deed to be detested than to be described. Blessed are they whose sins are covered, and happy are they that have learned to be ashamed. Unhappy are as many as cannot blush, but set such a face upon the matter that they glory in their sins and wickedness.

\index{Judgment, Last}After this, he touches briefly on the destruction of both the legates and those deceived by the legates, as well as those who fight against God and true religion to maintain Roman authority. The legates indeed assemble people of their faction for battle against the faithful, but the Lord has gathered them into a place called in Hebrew [illegible], which some interpret as the destruction of the River, and others as the army of desolation. But however that may be, the sense seems clear: they are assembled by the legates that they might withstand, or prohibit the destruction of the River and the ruin of Rome. But the Lord will also assemble the very same people, that in the same place and through the same means they may be destroyed by the Lord. This, we believe, will be accomplished finally at the last judgment. To the Lord Christ, our redeemer and revenger, be praise and glory. Amen.
